\documentclass[a4paper,10pt]{article}
\usepackage{geometry}
\geometry{top=0.55in,bottom=0.55in,left=0.7in,right=0.7in}
\usepackage{enumitem}
\usepackage[hidelinks]{hyperref}
\usepackage{fontspec}
\usepackage{xeCJK}
\usepackage{titlesec}
\usepackage{parskip}

\setmainfont{Latin Modern Roman}
\setCJKmainfont{FandolSong-Regular}[
  BoldFont=FandolSong-Bold,
  ItalicFont=FandolKai-Regular
]
\xeCJKsetup{CJKecglue=\hskip 0.10em plus 0.05em}

% Section: bold, large, with full-width rule beneath
\titleformat{\section}{\large\bfseries}{}{0em}{}[\vspace{-6pt}\rule{\linewidth}{0.5pt}]
\titlespacing*{\section}{0pt}{10pt}{4pt}

% Tight lists
\setlist[itemize]{leftmargin=1.5em, nosep, topsep=2pt, itemsep=1.5pt, parsep=0pt}

\setlength{\parindent}{0pt}
\setlength{\parskip}{0pt}
\pagestyle{empty}

% Helper: entry header with flush-right date
% Usage: \entry{Title}{Date}
\newcommand{\entry}[2]{%
  \noindent\textbf{#1}\hfill#2\par%
}
% Sub-line (advisor / subtitle), no date
\newcommand{\subline}[1]{%
  \noindent#1\par%
}

\begin{document}

%==================== Header ====================
\begin{center}
  {\LARGE\bfseries 赵宇骁}\\[5pt]
  Mail address: \href{mailto:zwb8@163.com}{zwb8@163.com}
  \quad\textbar\quad
  Mobilephone: 13316988235
\end{center}

\vspace{4pt}

%==================== 教育背景 ====================
\section*{教育背景}

\entry{香港中文大学}{09/2025 -- 至今}
专业：人工智能 \quad 学位：理学硕士\\
核心课程：人工智能基础、人工智能模型与应用、人工智能实践、生成式人工智能模型、
人工智能新兴科技、深度学习、强化学习、计算机视觉与模式识别、语音和语言处理

\vspace{5pt}

\entry{英国曼彻斯特大学}{09/2022 -- 06/2025}
专业：计算机科学与数学 \quad 学位：荣誉理学学士\\
核心课程：人工智能导论、编程导论、软件工程、图形与虚拟环境、概率与统计、
机器学习、随机过程、线性回归

\vspace{5pt}

\noindent\textbf{奖项与证书}
\begin{itemize}
  \item GreatUniHack 2023 编程竞赛参与证书，团队亚军 \hfill 11/2023
  \item 曼彻斯特大学 Stellify 卓越奖（前 5\% 学生） \hfill 06/2025
\end{itemize}

%==================== 项目经历 ====================
\section*{项目经历}

\noindent\textbf{硕士个人项目}

\vspace{3pt}
\entry{VLA-LocoMAN：面向四足机器人运动操作任务的端到端控制器学习}{11/2025 -- 至今}
\begin{itemize}
  \item 集成控制架构：构建一个分层强化学习框架，无缝协调行走与操作任务，使机器人能够在导航的同时执行复杂的物体操作，将 locomotion 与 manipulation 解耦并协同优化，实现移动中稳定抓取与操作。
  \item 稳健的仿真到现实迁移流程：开发迁移学习方法，使在仿真环境中训练的策略能够部署到实体机器人上，且性能损失最小，让算法在现实世界中发挥作用。
  \item 使用 PPO/SAC 等算法训练策略网络，结合 reward shaping, curriculum learning 改善训练稳定性与收敛效率。
  \item 视觉-语言动作接口：设计一个直观的端到端控制系统，用户可以通过结合视觉上下文的自然语言指令来操控机器人——让机器人控制人人可用，实现自然语言指令到连续控制动作的端到端映射。
\end{itemize}

\vspace{5pt}
\entry{生成式人工智能模型项目：大模型对齐与生成优化}{01/2026 -- 至今}
\begin{itemize}
  \item 深入研究 Transformer、Diffusion、LLM 对齐机制并实现 LoRA 微调与指令微调（Instruction Tuning）。
  \item 探索 RLHF / DPO 在模型对齐中的效果。
  \item 实验对比 Prompt Engineering、Chain-of-Thought、Self-Consistency 对推理性能的提升。
  \item 分析大模型 hallucination 成因与缓解策略。
\end{itemize}

\vspace{5pt}
\entry{AI 大模型实践项目}{09/2024 -- 12/2024}
\begin{itemize}
  \item 系统运行从基础数学实现到模型部署的完整 AI 工程流程，涵盖矩阵运算与卷积实现、深度学习框架构建。
  \item 运用所学知识进行模型训练与优化，设计 CNN 神经网络与 Attention 网络构建，将模型部署在服务器及 LLM 中，并使用 NLP 知识训练高性能、高准确率的 chatbox。
  \item 运用 GPU 编程（CUDA/Triton）、模型部署与数据管道构建，形成从算法开发到工程落地的端到端实践能力。
\end{itemize}

\vspace{5pt}
\entry{LuminaAI：基于 Transformer 的智能邮件分类与 AI 辅助收件箱系统}{09/2025 -- 12/2025}
\begin{itemize}
  \item 开发端到端智能邮件管理系统，集成 Gmail API、DistilBERT 分类模型及 LLM 驱动的动作建议与问答模块。
  \item 负责模型训练与实验分析：基于 12,000 条多类别邮件数据微调 DistilBERT，实现 99.67\% 的测试准确率与 99.67\% 的宏 F1 分数，显著优于传统方法。
  \item 设计两阶段 LLM 动作引擎（阅读层 + 动作层），实现邮件摘要、优先级矩阵分类及可操作建议（如回复、创建待办）的自动生成。项目代码已开源：\href{https://github.com/BoatHermit/email_ai_agent}{github.com/BoatHermit/email\_ai\_agent}，提供 Docker 化部署方案。
  \item 实现混合检索式问答系统，支持自然语言查询并返回带引用的答案，提升透明度和可追溯性。
\end{itemize}

\vspace{6pt}
\noindent\textbf{本科个人项目}

\vspace{3pt}
\entry{多传感器火灾检测与视觉自主导航机器人项目}{09/2024 -- 06/2025}
\begin{itemize}
  \item 使用 ROS Noetic 设计并实现自主火灾检测机器人，集成 YOLOv5 视觉识别、RGB/热成像融合、LiDAR 测距与 SLAM（GMapping）空间建图。
  \item 开发基于视觉的环境感知系统，利用计算机视觉与图像处理算法实现实时场景理解与决策。
  \item 结合感知与运动控制，通过动态代价地图与 A*/Dijkstra 算法实现低能见度环境下的危险感知导航。
  \item 优化视觉至导航延迟，提升系统响应速度；在 Gazebo 仿真中实现 93\% 火焰检测 F1 值与 92\% 闭环精度。
\end{itemize}

\vspace{5pt}
\entry{机器学习项目：空气质量分析}{02/2024 -- 06/2024}
\begin{itemize}
  \item 处理 3,304 条多传感器环境数据，采用缺失值填补与数据划分策略，构建适用于交叉验证的标准化数据集。
  \item 基于 NumPy 实现自定义二分类器，采用铰链损失与 L2 正则化，利用向量化运算提升计算效率。
  \item 构建并调优多层感知机（MLP）模型，以均方误差（MSE）为优化目标，探索不同网络结构对拟合性能的影响。
  \item 对比 SGD 与 Adam 优化器在MLP 上的训练动态，综合 MSE 与 $R^2$ 指标评估模型效果,选取最优超参数组合。
\end{itemize}

%==================== 研究经历 ====================
\section*{研究经历}

\entry{面向特征偏移客户端的自适应联邦归一化方法}{09/2025 -- 12/2025}
\subline{导师：LAU Wing Cheong 劉永昌教授，香港中文大学}
\begin{itemize}
  \item 独立提出 Adaptive Federated Normalization (AFN) 算法，针对联邦学习中客户端特征分布偏移问题，在局部与全局批归一化统计量间引入可调插值参数 $\alpha$，通信开销可忽略。
  \item 在含多种特征偏移的 5 客户端 MNIST 场景中评估：AFN 达 98.76\% 平均准确率，优于 FedAvg（97.84\%）和 FedBN（98.57\%）；统计显著性通过 $t$ 检验验证（$p < 0.05$）。
  \item 使用 PyTorch 实现完整联邦训练流程，代码已开源；论文《Adaptive Federated Normalization for Feature-Shifted Clients》为预印本（第一作者，2025）。
\end{itemize}

\vspace{5pt}
\entry{智慧城市：AI 驱动的城市交通与基础设施优化}{09/2025 -- 12/2025}
\subline{导师：林晓俊教授，香港中文大学}
\begin{itemize}
  \item 作为核心成员参与团队项目（Team \#27），系统梳理 AI 在智慧城市中的应用框架，涵盖智能交通信号控制（Alibaba City Brain，缩短 15.3\% 旅行时间）、基于强化学习的能源分配及数字孪生基础设施预测性维护。
  \item 分析新加坡、深圳、哥本哈根等城市实际部署案例，评估 AI 系统在通行效率、电网稳定性与公共交通准点率方面的量化效果；独立撰写报告核心章节，完成 20+ 篇文献调研与政策建议。
\end{itemize}

\vspace{5pt}
\entry{基于扩散模型的图像伪造定位研究}{07/2024 -- 08/2024}
\subline{导师：倪江群教授，广东中山大学}
\begin{itemize}
  \item 使用 Python 和 PyTorch 开发了用于高精度图像伪造检测的深度学习模型，重点关注像素级检测精度。
  \item 采用 \texttt{torch.randn()} 进行高斯噪声初始化，这是基于扩散模型的图像处理中的关键步骤。
  \item 使用 PyTorch 构建 U-Net 网络，集成卷积层、ReLU 激活函数和池化操作，用于特征提取与噪声抑制。
  \item 通过 PyTorch 张量运算利用 GPU 并行处理优化去噪任务，达成了高计算效率。
\end{itemize}

%==================== 实习经历 ====================
\section*{实习经历}

\entry{北京字节跳动科技有限公司}{06/2023 -- 07/2023}
\subline{数据分析实习生}
\begin{itemize}
  \item 使用 Pandas、NumPy、Matplotlib 构建数据可视化仪表盘，对抖音商城用户行为、画像特征与销售趋势进行多维度分析，为年度报告输出关键可视化成果。
  \item 设计并实施 A/B 测试，运用独立样本 $t$ 检验（$p < 0.05$）评估首页改版与个性化推荐算法对 CTR 和 CVR 的影响；利用卡方检验与多元线性回归量化用户活跃度与 GMV 的关联，定位核心驱动因素。
  \item 分析成果被采纳后，抖音商城用户次日留存率提升约 5\%。
\end{itemize}

%==================== 在校经历 ====================
\section*{在校经历}

\entry{香港中文大学人工智能专业课程代表}{09/2025 -- 至今}
\begin{itemize}
  \item 组织 AI 主题社交活动与工作坊，促进学生合作与项目参与。
  \item 撰写学术与社区提案，提升项目参与度与学生体验。
  \item 策划并发起 CUHK AI 标志设计竞赛，增强项目识别度与创新氛围。
\end{itemize}

\vspace{5pt}
\entry{曼彻斯特大学计算机与数学专业学生代表}{09/2023 -- 06/2025}
\begin{itemize}
  \item 代表 750+ 学生参加院系会议，传达学术反馈与教学建议。
  \item 提出评分改革方案并改进助教评分一致性, 发布每周通讯并通过问卷收集学生意见。
\end{itemize}

\vspace{5pt}
\entry{Stellify卓越奖}{06/2025}
\begin{itemize}
  \item 因领导力与社会责任获得 Stellify 卓越奖；完成 Grand Ethical Challenge 项目，并用大量时间投身社区服务与学生辅导志愿活动。
\end{itemize}

\vspace{5pt}
\entry{GreatUniHack 2023 编程竞赛}{09/2023 -- 11/2023}
\begin{itemize}
  \item 团队亚军。开发基于 Python 的人工智能 Web 应用，用于森林砍伐趋势检测与可视化；实现数据处理、图像标准化及后续的 CNN 处理，以及 Google Colab 云端部署。
\end{itemize}

%==================== 技术技能 ====================
\section*{技术技能}

\begin{itemize}
  \item \textbf{编程语言：} Python、Java、R、HTML/CSS
  \item \textbf{深度学习框架：} PyTorch、TensorFlow、scikit-learn、HuggingFace Transformers
  \item \textbf{机器学习与数据科学：} Pandas、NumPy、Matplotlib、TensorBoard
  \item \textbf{工程工具：} Git、Docker、CUDA/Triton、ROS Noetic、Gazebo、Jupyter Notebook、Google Colab
\end{itemize}

\end{document}
